\section{Machine model and reduction statement}\label{sec:model}

An input polynomial of degree less than $n$ is its $n$ coefficients in
ascending degree order, including leading zero coefficients. The output
has $2n-1$ coefficients. Time counts steps of one deterministic Turing
machine with a fixed finite alphabet and a fixed finite number of
one-dimensional tapes. The machine, its tapes and its finite tables do
not depend on $n$. All bit movement, setup and cleanup are included.
The number of fixed tapes and the constants may be extremely large.
This paper makes no claim for a specified two-tape machine or for
practical execution time.

Internal array interfaces receive a self-delimiting binary description
of their shape on descriptor tapes. In particular the tensor interface
receives $t,D$, and a batched evaluator receives $k,D,B_0$; parameter
processing and counter initialization are included in the bounds.
These are shape parameters computed by the calling algorithm, with no
input-dependent advice. At the top, an initial scan determines the raw
input length in $O(n)$ steps. A shape word of length $O(\log(D+1))$
does not have unit cost; its handling is charged to the nonempty
$D$-bit records as specified in the subroutines.

For a binary linear map $L$, $L^{\T}$ always denotes the transpose in
the displayed coefficient coordinates. Field-trace adjoints will be
converted to this transpose by explicit Gram maps. No generic
circuit-transposition or inverse-as-adjoint simulation is assumed.

\begin{contract}[Equal-chunk interchange]\label{con:swap}
For a fixed $0<\tau<1$, the permutation
\[
 [P]\times[2^u]\times[G]\times[2^u]\times[B]
 \longrightarrow [P]\times[2^u]\times[G]\times[2^u]\times[B],
 \quad(p,x,g,y,z)\longmapsto(p,y,g,x,z)
\]
has time $O(Vu^\tau)$, where $V=PGB2^{2u}$. The constant is uniform
in all positive dimensions, including $B=1$, and includes processing
the shape, temporary padding, unpadding and cleanup.
\end{contract}
This is the statement imported from Lemma 12 of \cite{OpenAI26}.
It is an external theorem dependency, not a conjecture about primes or
about polynomial multiplication. Section~\ref{sec:tensor} explains its
separate use in the tensor construction.

\begin{contract}[Uniform binary tensor transform]\label{con:tensor}
For some fixed $\beta<1$, put
\[
 H=\begin{pmatrix}1&1\\1&0\end{pmatrix}.
\]
For every $t\geq1$ and binary payload width $D\geq1$, one fixed machine
computes $H^{\otimes t}\otimes I_D$ in $O(2^tDt^\beta)$ time,
including all auxiliary representations, padding and cleanup. The
constant is independent of $D$.
\end{contract}
Since $H=H^{\T}$ and $H^3=I$, this same algorithm computes the binary
transpose, and two applications compute the inverse. If
$S=\left(\begin{smallmatrix}1&1\\0&1\end{smallmatrix}\right)$ and
$J=\left(\begin{smallmatrix}0&1\\1&0\end{smallmatrix}\right)$,
then $S=HJ$. Lucas' binomial congruence gives the Pascal matrix
\[
 (\mathcal P_{2^t})_{s,j}=\binom js\bmod2
       = (S^{\otimes t})_{s,j}.
\]
Thus one call to Contract~\ref{con:tensor} and reversal of the record
order compute $\mathcal P_{2^t}$. Its transpose places that reversal on
the other side. The payload can be a list of coefficient words: it
need not be one element of an enormous field.

\begin{lemma}[Rectangular reorientation]\label{lem:rectangle}
Contract~\ref{con:swap} permits the address-field rotation
$[2^a]\times[2^b]\to[2^b]\times[2^a]$ in
$O(V(a+b)^\gamma)$ time, where $\gamma=1/(2-\tau)<1$.
The bound includes an arbitrary opaque suffix payload.
\end{lemma}
\begin{proof}
Use the Euclidean block-rotation schedule. At widths $a\geq b$,
write $a=qb+r$, exchange the right field past the $q$ complete
$b$-bit chunks, then rotate the remaining pair of widths $b,r$.
Stop when the smaller positive width is less than $z$; move its
individual bits across the other field, preserving their order.
Single-bit moves cost $O(V)$ by a sequential split and merge.
The Euclidean remainders satisfy
\[
 \sum_i q_i r_i^\tau
 \leq z^{\tau-1}\sum_i q_i r_i
 \leq(a+b)z^{\tau-1}.
\]
The last sum telescopes, since $q_i r_i=r_{i-1}-r_{i+1}$.
The final individual-bit moves cost $O(Vz)$. Choose $z$ within a
constant factor of $(a+b)^\gamma$ using a power of two and integer
arithmetic on $\lceil\log_2(a+b)\rceil$. Both terms have the asserted
bound. Reverse the actual schedule for the transpose permutation.
\end{proof}

\begin{theorem}[Exact reduction]\label{thm:reduction}
Contracts~\ref{con:swap} and~\ref{con:tensor} imply an exact
$O(n\log n)$ algorithm for multiplying two degree-$<n$ polynomials
over $\F_2$ in the stated model. For either fixed operand, the ordinary
binary transpose of the multiplication map has the same bound.
\end{theorem}
Sections~\ref{sec:fields}--\ref{sec:decoding} prove this implication.
The removal of Contract~\ref{con:tensor} as a separate premise
is isolated in Section~\ref{sec:tensor} and Appendix~\ref{app:word}.
Theorem~\ref{thm:reduction} does not assume an arithmetic-operation
model with unit-cost extension-field scalars.

